\documentclass[10pt]{extarticle}
\usepackage[a4paper,margin=1.5cm]{geometry}
\usepackage[T1]{fontenc}
\usepackage{amsmath,amssymb}
\usepackage{graphicx}
\usepackage[export]{adjustbox}
\usepackage{multicol}
\usepackage{enumitem}
\usepackage{booktabs,array,tabularx}
\setlength{\columnsep}{0.7cm}
\setlength{\parindent}{0pt}
\setlength{\parskip}{3pt}
\pagestyle{empty}
\makeatletter
% figure registry (figdefs.tex is generated by the extract step)
\newcommand{\qfigdef}[4]{\@namedef{qf@file@#1}{#2}\@namedef{qf@w@#1}{#3}\@namedef{qf@wide@#1}{#4}}
\newcommand{\qf@check}[1]{\@ifundefined{qf@file@#1}{\PackageError{score2}{Unknown figure id #1}{}}{}}
\newcommand{\qf@img}[4][]{\includegraphics[width=#2\dimexpr\@nameuse{qf@w@#4}pt\relax,max width=#3,max height=0.45\textheight,#1]{\@nameuse{qf@file@#4}}}
% Full-width questions: finalize wraps every question that contains a wide figure in \qfullbegin...\qfullend,
% so the WHOLE question (stem, figures, options) is set across the page in reading order, then the two columns
% resume. The question's head - its number, stem and first figure (\qfullsplit marks the end of the head) -
% is kept on one page, so a number is never left at a page foot with its figure overleaf; the rest may break
% across pages normally, so no big gaps are left. Without a \qfullsplit, the whole question is kept together
% if it is under 45% of a page. Never write these yourself.
\newif\ifqf@full
\newif\ifqf@boxed
\newsavebox{\qf@box}
\newcommand{\qf@open}{\global\qf@fulltrue\par\global\qf@boxedtrue
  \setbox\qf@box\vbox\bgroup\hsize\textwidth\linewidth\textwidth}
\newcommand{\qf@emit}[1]{\par\egroup\global\qf@boxedfalse
  \ifdim\dimexpr\ht\qf@box+\dp\qf@box\relax<#1\textheight\noindent\usebox{\qf@box}\par\else\unvbox\qf@box\fi}
\newcommand{\qfullsplit}{\ifqf@boxed\qf@emit{0.85}\fi}
\newcommand{\qf@close}{\ifqf@boxed\qf@emit{0.45}\fi\global\qf@fullfalse}
\newcommand{\qfullbegin}{\end{multicols}\qf@open}
\newcommand{\qfullend}{\qf@close\begin{multicols}{2}}
\newcommand{\qfullstart}{\qf@open}
\newcommand{\qfulldone}{\qf@close}
% \qfig[scale]{id}: block figure, centred, directly where it is written (size set by the tool).
% Fallback if a wide figure is not inside a full-width question: break out of the columns in place - no floats.
\newcommand{\qfig}[2][1]{\qf@check{#2}%
  \ifqf@full\qf@blk{#1}{#2}\else
  \ifnum\@nameuse{qf@wide@#2}=1 \expandafter\@firstoftwo\else\expandafter\@secondoftwo\fi
  {\end{multicols}\begin{center}\qf@img{#1}{\textwidth}{#2}\end{center}\begin{multicols}{2}}%
  {\qf@blk{#1}{#2}}\fi}
\newcommand{\qf@blk}[2]{\par\smallskip{\centering\qf@img{#1}{\linewidth}{#2}\par}\smallskip}
% \qfigi[scale]{id}: inline figure (inside option lists / tables / minipages); never breaks columns.
\newcommand{\qfigi}[2][1]{\qf@check{#2}\qf@img[valign=c]{#1}{\linewidth}{#2}}
\makeatother
% structure macros - lint reads these to check completeness and figure placement
\newcommand{\papersec}[1]{\par\bigskip{\centering\Large\bfseries PAPER-#1\par}\medskip}
\newcommand{\partsec}[2]{\par\medskip\noindent{\large\bfseries PART-#1: #2}\par\smallskip}
\newcommand{\sect}[1]{\par\smallskip\noindent\textbf{#1}\par}
\newcommand{\qn}[1]{\par\smallskip\noindent\textbf{#1.}~}
\newcommand{\paraq}[2]{\par\smallskip\noindent\textbf{Paragraph for Questions #1 and #2}\par}
\newenvironment{qoptions}{\begin{itemize}[leftmargin=1.8em,itemsep=1pt,topsep=2pt,label={}]}{\end{itemize}}
\newcommand{\opt}[2]{\item[(#1)] #2}
% figure registry (generated by score2_tools.py finalize)
\qfigdef{p1_phy_s1_q03_1}{Score Advanced - Test 23 (14 Apr 2024)-img/p1_phy_s1_q03_1.png}{213.9}{0}
\qfigdef{p1_phy_s2_q05_1}{Score Advanced - Test 23 (14 Apr 2024)-img/p1_phy_s2_q05_1.png}{50.4}{0}
\qfigdef{p1_phy_s2_q06_1}{Score Advanced - Test 23 (14 Apr 2024)-img/p1_phy_s2_q06_1.png}{208.0}{0}
\qfigdef{p1_phy_s2_q07_1}{Score Advanced - Test 23 (14 Apr 2024)-img/p1_phy_s2_q07_1.png}{95.1}{0}
\qfigdef{p1_phy_s2_q09_1}{Score Advanced - Test 23 (14 Apr 2024)-img/p1_phy_s2_q09_1.png}{105.3}{0}
\qfigdef{p1_phy_s2_q10_1}{Score Advanced - Test 23 (14 Apr 2024)-img/p1_phy_s2_q10_1.png}{116.5}{0}
\qfigdef{p1_phy_s3_para1_2_1}{Score Advanced - Test 23 (14 Apr 2024)-img/p1_phy_s3_para1_2_1.png}{161.3}{0}
\qfigdef{p1_phy_s3_para3_4_1}{Score Advanced - Test 23 (14 Apr 2024)-img/p1_phy_s3_para3_4_1.png}{187.6}{0}
\qfigdef{p1_phy_s3_para5_6_1}{Score Advanced - Test 23 (14 Apr 2024)-img/p1_phy_s3_para5_6_1.png}{105.5}{0}
\qfigdef{p1_phy_s4_q07_1}{Score Advanced - Test 23 (14 Apr 2024)-img/p1_phy_s4_q07_1.png}{113.4}{0}
\qfigdef{p1_chem_s1_q02_A}{Score Advanced - Test 23 (14 Apr 2024)-img/p1_chem_s1_q02_A.png}{101.1}{0}
\qfigdef{p1_chem_s1_q02_B}{Score Advanced - Test 23 (14 Apr 2024)-img/p1_chem_s1_q02_B.png}{114.4}{0}
\qfigdef{p1_chem_s1_q02_C}{Score Advanced - Test 23 (14 Apr 2024)-img/p1_chem_s1_q02_C.png}{131.1}{0}
\qfigdef{p1_chem_s1_q02_D}{Score Advanced - Test 23 (14 Apr 2024)-img/p1_chem_s1_q02_D.png}{135.2}{0}
\qfigdef{p1_chem_s2_q07_A}{Score Advanced - Test 23 (14 Apr 2024)-img/p1_chem_s2_q07_A.png}{59.3}{0}
\qfigdef{p1_chem_s2_q07_B}{Score Advanced - Test 23 (14 Apr 2024)-img/p1_chem_s2_q07_B.png}{59.6}{0}
\qfigdef{p1_chem_s2_q07_C}{Score Advanced - Test 23 (14 Apr 2024)-img/p1_chem_s2_q07_C.png}{57.0}{0}
\qfigdef{p1_chem_s2_q07_D}{Score Advanced - Test 23 (14 Apr 2024)-img/p1_chem_s2_q07_D.png}{58.2}{0}
\qfigdef{p1_chem_s2_q08_1}{Score Advanced - Test 23 (14 Apr 2024)-img/p1_chem_s2_q08_1.png}{152.0}{0}
\qfigdef{p1_chem_s3_para3_4_1}{Score Advanced - Test 23 (14 Apr 2024)-img/p1_chem_s3_para3_4_1.png}{207.7}{0}
\qfigdef{p1_chem_s3_para5_6_1}{Score Advanced - Test 23 (14 Apr 2024)-img/p1_chem_s3_para5_6_1.png}{131.5}{0}
\qfigdef{p2_phy_s1_q01_1}{Score Advanced - Test 23 (14 Apr 2024)-img/p2_phy_s1_q01_1.png}{128.8}{0}
\qfigdef{p2_phy_s1_q02_1}{Score Advanced - Test 23 (14 Apr 2024)-img/p2_phy_s1_q02_1.png}{71.5}{0}
\qfigdef{p2_phy_s1_q03_1}{Score Advanced - Test 23 (14 Apr 2024)-img/p2_phy_s1_q03_1.png}{86.2}{0}
\qfigdef{p2_phy_s1_q04_1}{Score Advanced - Test 23 (14 Apr 2024)-img/p2_phy_s1_q04_1.png}{187.3}{0}
\qfigdef{p2_phy_s1_q06_1}{Score Advanced - Test 23 (14 Apr 2024)-img/p2_phy_s1_q06_1.png}{108.3}{0}
\qfigdef{p2_phy_s2_para7_8_1}{Score Advanced - Test 23 (14 Apr 2024)-img/p2_phy_s2_para7_8_1.png}{99.4}{0}
\qfigdef{p2_phy_s2_para9_10_1}{Score Advanced - Test 23 (14 Apr 2024)-img/p2_phy_s2_para9_10_1.png}{242.9}{0}
\qfigdef{p2_phy_s3_para1_2_1}{Score Advanced - Test 23 (14 Apr 2024)-img/p2_phy_s3_para1_2_1.png}{149.6}{0}
\qfigdef{p2_phy_s3_para3_4_1}{Score Advanced - Test 23 (14 Apr 2024)-img/p2_phy_s3_para3_4_1.png}{153.3}{0}
\qfigdef{p2_phy_s4_q07_1}{Score Advanced - Test 23 (14 Apr 2024)-img/p2_phy_s4_q07_1.png}{116.8}{0}
\qfigdef{p2_phy_s4_q09_1}{Score Advanced - Test 23 (14 Apr 2024)-img/p2_phy_s4_q09_1.png}{146.2}{0}
\qfigdef{p2_chem_s1_q03_1}{Score Advanced - Test 23 (14 Apr 2024)-img/p2_chem_s1_q03_1.png}{221.3}{0}
\qfigdef{p2_chem_s1_q03_A}{Score Advanced - Test 23 (14 Apr 2024)-img/p2_chem_s1_q03_A.png}{106.3}{0}
\qfigdef{p2_chem_s1_q03_B}{Score Advanced - Test 23 (14 Apr 2024)-img/p2_chem_s1_q03_B.png}{106.5}{0}
\qfigdef{p2_chem_s1_q03_C}{Score Advanced - Test 23 (14 Apr 2024)-img/p2_chem_s1_q03_C.png}{106.5}{0}
\qfigdef{p2_chem_s1_q03_D}{Score Advanced - Test 23 (14 Apr 2024)-img/p2_chem_s1_q03_D.png}{106.5}{0}
\qfigdef{p2_chem_s1_q04_1}{Score Advanced - Test 23 (14 Apr 2024)-img/p2_chem_s1_q04_1.png}{56.8}{0}
\qfigdef{p2_chem_s1_q04_2}{Score Advanced - Test 23 (14 Apr 2024)-img/p2_chem_s1_q04_2.png}{78.1}{0}
\qfigdef{p2_chem_s1_q04_3}{Score Advanced - Test 23 (14 Apr 2024)-img/p2_chem_s1_q04_3.png}{35.0}{0}
\qfigdef{p2_chem_s1_q04_4}{Score Advanced - Test 23 (14 Apr 2024)-img/p2_chem_s1_q04_4.png}{33.5}{0}
\qfigdef{p2_chem_s2_para9_10_1}{Score Advanced - Test 23 (14 Apr 2024)-img/p2_chem_s2_para9_10_1.png}{155.4}{0}
\qfigdef{p2_chem_s2_q09_1}{Score Advanced - Test 23 (14 Apr 2024)-img/p2_chem_s2_q09_1.png}{203.0}{0}
\qfigdef{p2_chem_s2_q09_2}{Score Advanced - Test 23 (14 Apr 2024)-img/p2_chem_s2_q09_2.png}{123.0}{0}
\qfigdef{p2_chem_s2_q09_3}{Score Advanced - Test 23 (14 Apr 2024)-img/p2_chem_s2_q09_3.png}{105.3}{0}
\qfigdef{p2_chem_s2_q09_4}{Score Advanced - Test 23 (14 Apr 2024)-img/p2_chem_s2_q09_4.png}{114.1}{0}
\qfigdef{p2_chem_s3_para3_4_1}{Score Advanced - Test 23 (14 Apr 2024)-img/p2_chem_s3_para3_4_1.png}{233.9}{0}
\qfigdef{p2_chem_s4_q09_1}{Score Advanced - Test 23 (14 Apr 2024)-img/p2_chem_s4_q09_1.png}{316.2}{1}
\qfigdef{p2_math_s2_para9_10_1}{Score Advanced - Test 23 (14 Apr 2024)-img/p2_math_s2_para9_10_1.png}{161.3}{0}
\raggedcolumns
\makeatletter\@beginparpenalty=10000\makeatother
\newcolumntype{L}[1]{>{\raggedright\arraybackslash}p{#1}}
\begin{document}
\begin{center}{\Large\bfseries Score Advanced --- Test 23 (14 Apr 2024)}\end{center}
\begin{multicols}{2}
\raggedcolumns
\papersec{1}
\partsec{1}{PHYSICS}
\sect{SECTION-I (i)}
\qn{1} Two conductor are embedded in a material of conductivity $\sigma$ and dielectric constant $\varepsilon_r$. The resistance between the two conductors is measured to be R. The capacitance between two conductors is :
\begin{qoptions}
\opt{A}{$\dfrac{\varepsilon_0\varepsilon_r}{2\sigma R}$}
\opt{B}{$\dfrac{2\varepsilon_0\varepsilon_r}{\sigma R}$}
\opt{C}{$\dfrac{\varepsilon_0\varepsilon_r}{\sigma R}$}
\opt{D}{$\dfrac{5\varepsilon_0\varepsilon_r}{2\sigma R}$}
\end{qoptions}
\qn{2} At a certain temperature T, a liquid with volume $V_1$ and volumetric expansion coefficient $\alpha$ is at a certain level in a container. At the same temperature, a solid body with a volume $V_2$ and coefficient of linear expansion $\lambda$ remains outside the liquid. The temperature of the liquid in the container is reduced by $\Delta$T, and the temperature of the object is raised by an equal amount and then thrown into the container. In this case, the liquid level in the container does not change. How much $\Delta$T is required to achieve this situation? The densities and volumes of the object and the liquid are such that they allow the object to fully sink into the liquid. It will be assumed that the container does not expand, is heat insulating and that there will be no heat exchange between the liquid and the object.
\begin{qoptions}
\opt{A}{$\dfrac{V_1-V_2}{\alpha V_1-3\lambda V_2}$}
\opt{B}{$\dfrac{V_1-V_2}{\alpha V_1-\lambda V_2}$}
\opt{C}{$\dfrac{V_1}{\alpha V_1-3\lambda V_2}$}
\opt{D}{$\dfrac{V_2}{\alpha V_1-3\lambda V_2}$}
\end{qoptions}
\qn{3} In the figures below, images of the two objects on the left, created by a converging lens, are given on the right. Assuming that the image given as K is correct, which of the following options is correct?\par
\qfig{p1_phy_s1_q03_1}
\begin{qoptions}
\opt{A}{L, M and N}
\opt{B}{L and M}
\opt{C}{L and O}
\opt{D}{M, N and O}
\end{qoptions}
\qn{4} At normal incidence, a beam of light propagating in vacuum reflects off an interface with a medium of refractive index n = 2.0. The fraction of energy reflected R is given by $R=\left(\dfrac{n-1}{n+1}\right)^2$. If the percentage error in the value of n is 1\%, the percentage error in the estimation of R is :-
\begin{qoptions}
\opt{A}{6\%}
\opt{B}{$\dfrac{16}{3}$ \%}
\opt{C}{4\%}
\opt{D}{$\dfrac83$\%}
\end{qoptions}
\sect{SECTION-I (ii)}
\qn{5} Sound wave is incident on Air water interface. If equation of incident disturbance at interface is S = $S_0$ sin ($\omega$t). Here V and B are speed of sound and Bulk modulus respectively.\par
\qfig{p1_phy_s2_q05_1}
\begin{qoptions}
\opt{A}{If $2B_2V_1=B_1V_2$, amplitude of reflected disturbance will be $\dfrac{S_0}{3}$.}
\opt{B}{If $2B_2V_1=B_1V_2$, amplitude of transmitted disturbance will be $\dfrac{4S_0}{3}$.}
\opt{C}{If $B_2V_1=2B_1V_2$, amplitude of reflected disturbance will be $\dfrac{S_0}{3}$.}
\opt{D}{If $B_2V_1=2B_1V_2$, then amplitude of transmitted disturbance will be $\dfrac{2S_0}{3}$}
\end{qoptions}
\qn{6} We have two piston cylinder arrangement connected by a valve as shown in figure.\par
\qfig{p1_phy_s2_q06_1}
Same ideal monoatomic gas is filled in both cylinders at same temperature $T_0$ but at different pressures $2P_0$ and $P_0$. Both cylinders are connected by a thermal tube fitted with a valve, which is initially closed. Now valve is opened and gas starts leaking from cylinder A to B. Mr. Swapnil keeps pushing cylinder A and maintain $2P_0$ pressure in cylinder A while Mr. Sri holds the piston B at its original position during the motion. Process is stopped when pressure in both cylinders becomes $2P_0$. (Neglect any loss of heat in surrounding)
\begin{qoptions}
\opt{A}{Final temperature of gas in both cylinders is $\dfrac{19}{15}T_0$.}
\opt{B}{Final volume of cylinder A is $\dfrac{3}{10}$V.}
\opt{C}{Final temperature of gas in both cylinders $\dfrac{17}{15}T_0$.}
\opt{D}{Final volume of cylinder A is $\dfrac{7}{10}V_0$.}
\end{qoptions}
\qn{7} In the parallel circuit shown in figure, L = 0.2 H and C = 20 $\mu$F. Choose the correct options.\par
\qfig{p1_phy_s2_q07_1}
\begin{qoptions}
\opt{A}{If R = 0 $\Omega$, resonant angular frequency is 500 rad/s.}
\opt{B}{If R = 50 $\Omega$, resonant angular frequency is 433 rad/s.}
\opt{C}{In above circuit resonance is not possible for any value of R.}
\opt{D}{At resonance, voltage and current in the circuit are in same phase.}
\end{qoptions}
\qn{8} Angular width of central maximum in the Fraunhofer diffraction pattern of a slit is measured. The slit is illuminated by light of wavelength 6000\AA.\par
Case-1 : When the slit is illuminated by light of another wavelength, the angular width decreases by 30\%.\par
Case-2 : When the original apparatus is immersed in a liquid, angular width decreases by 30\%.\par
Choose the correct option(s) :
\begin{qoptions}
\opt{A}{In case 1, wavelength of light is 4800\AA.}
\opt{B}{In case 1, wavelength of light is 4200\AA.}
\opt{C}{In case 2, refractive index of liquid is 1.25.}
\opt{D}{In case 2, refractive index of liquid is 1.43.}
\end{qoptions}
\qn{9} If the capacitance ratio $\dfrac{C_2}{C_1}=\eta$ then the potential difference V between point M and N in the given figure.\par
\qfig{p1_phy_s2_q09_1}
\begin{qoptions}
\opt{A}{If E = 110 V and $\eta$ = 2 then V = 10 volt}
\opt{B}{If E = 285V and $\eta$ = 3 then V = 15 volt}
\opt{C}{If E = 110 V and $\eta=\dfrac12$ then V = 10 volt}
\opt{D}{If E = 380 V and $\eta$ = 3 then V = 20 volt.}
\end{qoptions}
\qn{10} In order to further increase the energy of the cyclotron, scientists built a ``sector- shaped focusing cyclotron''. During the sector-shaped focusing process, ions can rotate periodically in a closed equilibrium orbit at a constant speed.\par
A simplified diagram of the fan-shaped focusing magnetic field distribution is shown in the figure. The circular area with the mouth at the center is equally divided into six fan- shaped areas, three of which are peak areas and three are valley areas. The peak areas and valley areas are alternately distributed. In the peak area, there is a uniform magnetic field with a direction perpendicular to the plane of the paper. The magnetic induction intensity is B. There is no magnetic field in the valley area. Positive ions with mass m and charge q rotate at a constant speed v. Their closed equilibrium orbit is shown by the dotted line in the figure.\par
\qfig{p1_phy_s2_q10_1}
\begin{qoptions}
\opt{A}{The radius of orbit in closed equilibrium orbit in peak area is $\dfrac{mv}{qB}$.}
\opt{B}{The time period of ions rotating around the closed equilibrium orbit is $\dfrac{(2\pi+3\sqrt3)m}{qB}$.}
\opt{C}{The centre angle of the arc ($\theta$) of the orbit in peak area is $\dfrac\pi3$.}
\opt{D}{Now a uniform magnetic field perpendicular to the paper plane is also applied in the valley area. The magnetic induction intensity is B$'$. The central angle of the new closed equilibrium orbit (as found in option (C)) in a peak area is now 90$^\circ$. The ratio $\dfrac{B'}{B}$ is $\dfrac{\sqrt3-1}{2}$.}
\end{qoptions}
\sect{SECTION-II (i)}
\paraq{1}{2}
A square single-turn loop carrying a current i lies in (x, y) plane with its side on the wires x = $\pm$0.5 m and y = $\pm$0.5 m. A small 5 turn search coil of effective area 0.05 m$^2$ is placed in xy plane at the origin.\par
\qfig{p1_phy_s3_para1_2_1}
\qn{1} Calculate the maximum instantaneous emf induced in search coil (in mV) if current (i) in loop is given as i = 20 sin(1000 t).
\qn{2} Calculate the rms value of the search coil emf (in mV) if i = 20 sin (1000 t) + 5 sin (3000 t).
\paraq{3}{4}
Fig shows a cylindrical body floating on a liquid surface. Line AB is passing through diameter of circular part of cylinder. Assume liquid and material of body makes a non-wetting combination with contact angle 143$^\circ$ at the circumference of cylinder passing through line AB. The line AB is at a depth h below the horizontal flat part of free surface. Neglect the effect of surface tension at the flat ends of cylinder. [surface tension = 0.05 N/m, density of liquid = 4 $\times$ 10$^3$ kg/m$^3$, radius of cylinder = 2mm, g = 10 m/s$^2$]\par
\qfig{p1_phy_s3_para3_4_1}
\qn{3} Find value of h (in mm) as shown in figure.
\qn{4} Find density of material of cylinder in (kg/m$^3$)
\paraq{5}{6}
An expedition team member encountered a ravine during the expedition. One side of the ravine was vertical and the slope on the other side was parabolic. The team member started from the vertical side of the ravine with a speed $v_0$ in the horizontal direction and jumps to the slope on the other side. As shown in the figure, the coordinate system Oxy is established with the O point at the bottom of the ravine as the origin. It is known that the height of the vertical side of the ravine is 2h, and the parabola equation of the slope is $y=\dfrac{x^2}{2h}$\par
The mass of the expedition team member is m. The human being is regarded as a particle, the air resistance is ignored, and the acceleration due to gravity is g. (Take h = 10 m, m = 1 kg)\par
\qfig{p1_phy_s3_para5_6_1}
\qn{5} At what speed (in m/s) does this person jump out horizontally, his kinetic energy when he lands on the slope is minimum?
\qn{6} What is the minimum value of kinetic energy (in J) when he lands on the slope?
\sect{SECTION-II (ii)}
\qn{7} A horizontal ray of light passes through a prism of refractive index 2.0 and apex angle 6$^\circ$ and then strikes a vertical mirror as shown in the figure. Through what angle (in degree) must the mirror be rotated if after reflection the ray is to be horizontal?\par
\qfig{p1_phy_s4_q07_1}
\qn{8} A mass of 1 kg is suspended from a spring of spring constant 25 Nm$^{-1}$. If the undamped (or natural) frequency is $\dfrac{2}{\sqrt3}$ times the damped frequency, calculate the damping constant in kg s$^{-1}$.
\qn{9} A piece of lead of mass 0.600 kg and at a temperature of 373.0 K is placed in a copper vessel of mass 0.200 kg which contains of 0.500 kg of water at a temperature of 290.0 K. The final steady temperature of the mixture is 292.8 K. The specific heat capacity of copper is 386 J kg$^{-1}$K$^{-1}$ and that the water is 4180 J kg$^{-1}$K$^{-1}$. Calculate the specific heat capacity of lead (in J kg K$^{-1}$). Assume that there is no work interactions and there is no heat exchange with the surrounding. [Give answer to nearest integer]
\partsec{2}{CHEMISTRY}
\sect{SECTION-I (i)}
\qn{1} Select the correct statements\par
(I) 1 gm - molecule of any gas contains equal number of molecules\par
(II) Atomic mass of carbon-12 is 12u\par
(III) During a reaction, moles and mass of atoms remain constant\par
(IV) Both molality and mole fraction are temperature independent
\begin{qoptions}
\opt{A}{only I, II, III}
\opt{B}{only III, IV}
\opt{C}{only I, III, IV}
\opt{D}{I, II, III and IV}
\end{qoptions}
\qn{2} Identify the reaction(s) which yield diastereomers as a major product?
\begin{qoptions}
\opt{A}{\qfigi{p1_chem_s1_q02_A}}
\opt{B}{\qfigi{p1_chem_s1_q02_B}}
\opt{C}{\qfigi{p1_chem_s1_q02_C}}
\opt{D}{\qfigi{p1_chem_s1_q02_D}}
\end{qoptions}
\qn{3} Among the following how many of given is/are not used in preparation of NH$_3$ through Haber's Process\par
N$_2$, H$_2$, K$_2$O, Al$_2$O$_3$, CaF$_2$, KCl, FeCO$_3$
\begin{qoptions}
\opt{A}{2}
\opt{B}{3}
\opt{C}{6}
\opt{D}{5}
\end{qoptions}
\qn{4} Find out number of square planer species among the following.\par
$[\text{ICl}_4]^{-}$, SF$_4$, cis $[\text{PtCl}_2(\text{PMe}_3)_2]$, trans $[\text{PtCl}_2(\text{PMe}_3)_2]$, SO$_2$Cl$_2$, SNF$_3$
\begin{qoptions}
\opt{A}{2}
\opt{B}{3}
\opt{C}{4}
\opt{D}{5}
\end{qoptions}
\sect{SECTION-I (ii)}
\qn{5} Which of the following statements is/are correct
\begin{qoptions}
\opt{A}{All adiabatic process are iso-entropic process}
\opt{B}{All spontaneous process must be exothermic in nature}
\opt{C}{Reversible isothermal process for an ideal gas is iso-enthalpic process}
\opt{D}{Heat capacity of any substance undergoing adiabatic process is zero}
\end{qoptions}
\qn{6} If de-Brogle wavelength of 1$^{st}$ excited state of He$^{+}$ ion is X, then select the correct options
\begin{qoptions}
\opt{A}{Radius of ground state of hydrogen atom is X/2$\pi$}
\opt{B}{Radius of ground state of He$^{+}$ ion is X/4$\pi$}
\opt{C}{de-Brogle wavelength of 2$^{nd}$ excited of Li$^{+2}$ ion is X}
\opt{D}{Ratio of de-Broglie wavelength of 1$^{st}$ excited state of He$^{+}$ ion to 2$^{nd}$ excited state of Li$^{+2}$ is 1 : 1}
\end{qoptions}
\qn{7} On treatment with phenylhydrazine, aldohexoses \textbf{P} and \textbf{Q} give the same osazone. On treatment with warm nitric acid, \textbf{P} gives an optically inactive aldaric acid, but sugar \textbf{Q} gives an optically active aldaric acid. Sugars \textbf{P} and \textbf{Q} are both degraded to aldopentose \textbf{R}, which gives an optically active aldaric acid on treatment with nitric acid. Aldopentose \textbf{R} is degraded to aldotetrose \textbf{S}, which gives optically active tartaric acid when it is treated with nitric acid. Aldotetrose \textbf{S} is degraded to D(+)-glyceraldehyde. Identify correct structures of sugars ?
\begin{qoptions}
\opt{A}{\qfigi{p1_chem_s2_q07_A}}
\opt{B}{\qfigi{p1_chem_s2_q07_B}}
\opt{C}{\qfigi{p1_chem_s2_q07_C}}
\opt{D}{\qfigi{p1_chem_s2_q07_D}}
\end{qoptions}
\qn{8} \qfig{p1_chem_s2_q08_1}
Considering above compounds, identify the correct statement(s)?
\begin{qoptions}
\opt{A}{Compound (\textbf{P}) and Compound (\textbf{Q}) can be differentiated by NaHCO$_3$ as well as aq.NaOH}
\opt{B}{Compound (\textbf{Q}) and Compound (\textbf{R}) can be differentiated by Neutral FeCl$_3$}
\opt{C}{Compound (\textbf{R}) and Compound (\textbf{S}) can be differentiated by NaHSO$_3$}
\opt{D}{Only compound (\textbf{R}) gives Mustard oil smell with CHCl$_3$ -KOH}
\end{qoptions}
\qn{9} Correct statement/statements among the following is/are
\begin{qoptions}
\opt{A}{N$_2$F$_2$ \& have two geometrical forms and both are planar}
\opt{B}{BrF$_3$ is used as an aprotic solvent}
\opt{C}{Hydrolysis of K$_2$S$_2$O$_8$ with D$_2$O gives D$_2$O$_2$ \& KDSO$_4$.}
\opt{D}{Cu, Sc \& Hg are transition elements}
\end{qoptions}
\qn{10} Acidified permanganate solution oxidises
\begin{qoptions}
\opt{A}{oxalates to carbon dioxide}
\opt{B}{Iodide to iodate}
\opt{C}{nitrite to nitrate}
\opt{D}{Fe$^{2+}$ ion to Fe$^{3+}$ ion}
\end{qoptions}
\sect{SECTION-II (i)}
\paraq{1}{2}
Salt hydrolysis : This is the reverse reaction of acid base neutralisation\par
Salt of many types :\par
(i) Salt of strong acid and strong base\par
(ii) Salt of strong base and weak acid\par
(iii) Salt of strong acid and weak base\par
(iv) Salt of weak acid and weak base\par
Answer the following questions
\qn{1} 0.2 M aq. solution of Na$_2$CO$_3$ is mixed with 0.2 M aq. HCl in equal volume proportion at 25$^\circ$C. pH of the resulting solution will be\par
(For H$_2$CO$_3$: K$_{a_1}$ = 2 $\times$ 10$^{-6}$, K$_{a_2}$ = 8 $\times$ 10$^{-10}$, log (2) = 0.3)
\qn{2} In the above final solution, if concentration of CO$_3^{2-}$ is X molar, then value of 10 X will be
\paraq{3}{4}
Reaction scheme:\par
\qfig{p1_chem_s3_para3_4_1}
\qn{3} Total number of bromine atoms present in product when compound(Q) react with Br$_2$ $-$ H$_2$O ?
\qn{4} Total number of Nitrogen atoms present in compound(U)?
\paraq{5}{6}
For the reaction 3A $\rightarrow$ 2B, the following graph is obtained experimentally.\par
\qfig{p1_chem_s3_para5_6_1}
\qn{5} What is order of reaction.
\qn{6} The rate of reaction when [A] = 0.2 M is P $\times$ 10$^{-4}$ M min$^{-1}$, the value of `9 P' is.
\sect{SECTION-II (ii)}
\qn{7} The primary alcohol 1-Propanol can be oxidized with aqueous nitric acid. By product of this redox reaction is nitric oxide. How much weight in grams of nitric acid is required for complete oxidation of 180 grams of 1-Propanol ? \textbf{(\% reaction yield is 80\%)}
\qn{8} Hard water can be softened by using different reagents, how many among the following are not used for softening of water.\par
(1) Hydrated sodium alumino silicates\par
(2) Acetanilides\par
(3) Sodium hexametaphosphates\par
(4) Washing soda\par
(5) lime\par
(6) A mixture of marble \& Acrolein
\qn{9} At a certain temperature for which RT = 30 lit atm mol$^{-1}$, the density of a real gas (in gm lit$^{-1}$) is:\par
d (in gm lit$^{-1}$) = 2P + 0.02 P$^2$\par
where P is the pressure in atm, then molar mass of the gas (in gm/mol) will be
\partsec{3}{MATHEMATICS}
\sect{SECTION-I(i)}
\qn{1} If $\alpha,\beta$ are roots of the equation $x^2-4x-3=0$, then find\par
\resizebox{\linewidth}{!}{$\dfrac{\alpha^{-\alpha\beta}\left[\alpha^{\alpha}\alpha^{\beta}-4\alpha^{-\alpha\beta}\right]+\beta^{-\alpha\beta}\left[\beta^{\alpha}\beta^{\beta}-4\beta^{-\alpha\beta}\right]}{\alpha.\alpha^{\alpha+\beta}+\beta.\beta^{\alpha+\beta}}$}
\begin{qoptions}
\opt{A}{3}
\opt{B}{4}
\opt{C}{5}
\opt{D}{7}
\end{qoptions}
\qn{2} In a sequence $\{a_n\}$, $a_1=2$, $a_{n+1}=1-\dfrac1{a_n}$ for n $\geqslant$ 1. Let $P_n$ be the product of its first n terms, then the value of $P_{2017}$ is
\begin{qoptions}
\opt{A}{$\dfrac{-1}{2}$}
\opt{B}{2}
\opt{C}{$\dfrac12$}
\opt{D}{1}
\end{qoptions}
\qn{3} If $\tan3^\circ+2\tan6^\circ+4\tan12^\circ+8\cot24^\circ=\cot\theta^\circ$, then
\begin{qoptions}
\opt{A}{$\cot(10\theta)^\circ=1$}
\opt{B}{$\cot(15\theta)^\circ=1$}
\opt{C}{$\cot\theta^\circ=0$}
\opt{D}{$\cot(15\theta)^\circ=\sqrt3$}
\end{qoptions}
\qn{4} A straight line through the point A($-$1, 3) cuts the line x + y = 4 and y = x at B and C respectively. If AB.AC = 16 and possible slopes of the line are $m_1$ \& $m_2$, where $-1\leqslant m_1,m_2\leqslant1$, then $\left(\dfrac{m_1-m_2}{1+m_1m_2}\right)^2$ is equal to
\begin{qoptions}
\opt{A}{1}
\opt{B}{$\dfrac34$}
\opt{C}{3}
\opt{D}{$\dfrac13$}
\end{qoptions}
\sect{SECTION-I(ii)}
\qn{5} From any point on the line y = x + 4 tangents are drawn to the auxillary circle of the ellipse $x^2+4y^2=4$. If P, Q are the points of contact \& A, B are the corresponding points of P \& Q on the ellipse respectively, then
\begin{qoptions}
\opt{A}{Locus of mid point of AB is ellipse with eccentricity $\dfrac12$}
\opt{B}{Locus of mid point of AB is ellipse with eccentricity $\dfrac{\sqrt3}{2}$}
\opt{C}{Locus of mid point of AB is ellipse $4x^2+y^2+x-2y=0$}
\opt{D}{Locus of mid point of AB is ellipse $x^2+4y^2+x-2y=0$}
\end{qoptions}
\qn{6} If tangents of conic $x^2+4xy+3y^2-5x-6y+3=0$ parallel to line x + 4y = 0 are $T_1$ \& $T_2$, then
\begin{qoptions}
\opt{A}{Distance between $T_1$ \& $T_2$ is $\dfrac2{\sqrt{17}}$}
\opt{B}{Distance between $T_1$ \& $T_2$ is $\dfrac3{\sqrt{17}}$}
\opt{C}{Area enclosed by $T_1$, $T_2$ \& coordinate axis is $\dfrac{39}8$}
\opt{D}{Area enclosed by $T_1$, $T_2$ \& coordinate axis is 5.}
\end{qoptions}
\qn{7} Let \par\resizebox{\linewidth}{!}{$I_n=\displaystyle\int_0^{\pi}\dfrac{2\cos x-\cos(2n-1)x-2\cos(2n)x-\cos(2n+1)x+2}{1-\cos2x}dx$}\par Then (n $\in$ W)
\begin{qoptions}
\opt{A}{$\displaystyle\sum_{n=0}^{99}I_n=9900\pi$}
\opt{B}{$I_{1011}=2022\pi$}
\opt{C}{$\displaystyle\sum_{n=1}^{100}I_n=10100\pi$}
\opt{D}{$I_{1011}=1011\pi$}
\end{qoptions}
\qn{8} If $\displaystyle\int\dfrac{e^{x-1}}{x^2-5x+4}2x\,dx=Af(x-1)+Bf(x-4)+C$ \& $f(x)=\displaystyle\int\dfrac{e^x}{x}dx$, then
\begin{qoptions}
\opt{A}{$A=-\dfrac23$}
\opt{B}{$B=\dfrac43e^3$}
\opt{C}{$A=\dfrac23$}
\opt{D}{$B=\dfrac83e^3$}
\end{qoptions}
\qn{9} Let f(x) satisfies differential equation $y=y'+y''+y'''+\dots\infty$, $f(0)=\dfrac1e$ and g(x) satisfies differential equation $y+y'=y''+y'''+\dots\infty$, then
\begin{qoptions}
\opt{A}{$g(x)=g''(x)$}
\opt{B}{$g(x)=2g''(x)$}
\opt{C}{f(x) = 0 has one root}
\opt{D}{$f(x)+x^3=0$ has one root}
\end{qoptions}
\qn{10} Possible value(s) of ``a'' if equation $1-\cos x=\dfrac{\sqrt3}{2}|x|+a$, x $\in$ (0, $\pi$), has exactly one solution is/are
\begin{qoptions}
\opt{A}{$\dfrac12-\dfrac{\pi}{2\sqrt3}$}
\opt{B}{$\dfrac32-\dfrac{\pi}{\sqrt3}$}
\opt{C}{$1-\dfrac{\sqrt3\pi}{4}$}
\opt{D}{$1-\dfrac1{\sqrt2}-\dfrac{\sqrt3\pi}{8}$}
\end{qoptions}
\sect{SECTION-II (i)}
\paraq{1}{2}
Consider a vector $\vec r=\vec a+\vec a\times\hat b$, where the projection of $\vec r$ on $\vec a$ is 2.\par
On the basis of above information, answer the following questions:
\qn{1} $|\vec a|$ is equal to
\qn{2} Maximum value of $|\vec r|^2$ is
\paraq{3}{4}
Consider two straight lines\par
$L_1:\vec r=\left(2\hat i+2\hat j+4\hat k\right)+\lambda\left(-\dfrac12\hat i+\dfrac12\hat j+\hat k\right)$ \&\par
$L_2:\vec r=\left(-\hat j+3\hat k\right)+\mu\left(\dfrac12\hat i+\hat j-\hat k\right)$\par
Let P and Q be the variable points on $L_1$ and $L_2$ respectively and point M divides PQ in the ratio 2 : 3 internally. Also R and S be the variable points on line $L_1$ and $L_2$ respectively which are nearest to each other.
\qn{3} The square of distance of origin from the line passing through ($-$1, 0, 2) and is parallel to the line joining R and S is
\qn{4} If the locus of M is 2x + z = k, then k is equal to :-
\paraq{5}{6}
Let z is a complex number satisfying $|z-3-4i|=4$ and $\theta_1$ and $\theta_2$ are maximum and minimum value of amplitude of z respectively. ($-\pi<\theta_1,\theta_2\leq\pi$)\par
On the basis of above information, answer the following questions :
\qn{5} The value of $|7\tan(\theta_1+\theta_2)|$ is
\qn{6} Maximum possible area of a triangle formed by z, iz \& z + iz is P, then value of [P], where [.] is G.I.F
\sect{SECTION-II (ii)}
\qn{7} Two circles whose equations are $x^2+y^2=9$ and $x^2+y^2-8x-6y+n^2=0$, n $\in$ I, have exactly two common tangents. The number of values of `n' is
\qn{8} If the number of solutions of $|3[x]-4x|=4$ is $\lambda$, then value of $\lambda$ is (where [.] is greatest integer function)
\qn{9} Given that P(A) = 2/7, P(B) = 2/7, $P\left(A\cap\bar B\cap C\right)=1/14$, $P\left(A\cap\bar B\cap\bar C\right)=1/14$. If $P\left(\bar A|\bar B\right)=\dfrac xy$ where x, y are coprime positive integers, then x + y =
\raggedcolumns
\papersec{2}
\partsec{1}{PHYSICS}
\sect{SECTION-I (i)}
\qn{1} A uniform rod PR of mass M and length L, area of cross-section A is placed on a smooth horizontal surface. The forces acting on the rod are shown in the figure. Here Y is Young's modulus of elasticity of the rod. Choose the correct options. (Given that PQ = QR = $\dfrac L2$)\par
\qfig{p2_phy_s1_q01_1}
\begin{qoptions}
\opt{A}{Ratio of elongation in section PQ of rod and section QR of rod is 5 : 7.}
\opt{B}{Elastic potential energy stored in the PQ section of rod is $\dfrac{19F^2L}{48AY}$.}
\opt{C}{Ratio of elastic potential energy stored in section PQ and section QR of the rod is 29 : 37.}
\opt{D}{The total elastic potential energy stored in the rod is $\dfrac{7F^2L}{6AY}$.}
\end{qoptions}
\qn{2} A fixed hollow frictionless cone is positioned with its tip pointing down. A small particle of mass m is released from rest on the inside surface from height 2H measured from vertex O of the cone. After it has slid half-way down to the tip, it bounces elastically off a platform. The platform is positioned at a 45$^\circ$ angle along the surface of the cone, so that the particle ends up being deflected horizontally along the surface of the cone (into the plane of the diagram shown). The particle then swirls up and around the cone before coming down moving in circular paths and always remaining in contact with the cone's inner surface.\par
\qfig{p2_phy_s1_q02_1}
\begin{qoptions}
\opt{A}{The particles maximum swirling height from point `O' after its collision with the platform is $\left(\dfrac{\sqrt5+1}{2}\right)$H.}
\opt{B}{The particles maximum swirling height from point `O' after its collision with the platform is $\dfrac{H}{\sqrt5+1}$.}
\opt{C}{The speed of the particle at the maximum height from point O after its collision with the platform is $\dfrac{\sqrt{8gH}}{\sqrt5+1}$.}
\opt{D}{The speed of the particle at the maximum height from point O after its collision with the platform is $\left(\sqrt5+1\right)\sqrt{8gH}$.}
\end{qoptions}
\qn{3} A uniform hemispherical bowl of mass `M' is at rest on frictionless horizontal ground surface. There is a small insect of mass `m' on bowl at point `A' of the bowl at rest. Now the insect moves on the hemisphere with constant speed `v' relative to bowl in vertical plane. Assume that insect does not slip on the bowl.\par
\qfig{p2_phy_s1_q03_1}
\begin{qoptions}
\opt{A}{When the insect reaches point `P' of the hemisphere the displacement of hemisphere w.r.t. ground is $\dfrac{mR}{(M+m)}(1+\cos\theta)$ horizontally towards left.}
\opt{B}{When the insect is at point `P' of the hemispherical bowl the acceleration of bowl is $\left(\dfrac{m}{M+m}\right)\dfrac{V^2\cos\theta}{R}$ horizontally towards left.}
\opt{C}{When the insect is at point `P' of the hemispherical bowl the acceleration of bowl is $\left(\dfrac{m}{M+m}\right)\dfrac{V^2\cos\theta}{2R}$ horizontally towards left.}
\opt{D}{When insect is at point `P' on the bowl the displacement of centre of mass of the system (bowl + insect) is $\dfrac{mR\sin\theta}{M+m}$ vertically upwards.}
\end{qoptions}
\qn{4} Water is supplied by a pump at a steady rate. Inlet and outlet pipes of common access attached to the pump are horizontal and have diameters of cross-sections 2d and d respectively as shown in the figure. Assume that flow through the entries system is steady and that the viscous effects are neglected. Velocity of water at the inlet is u = 1m/s and pump supplies mechanical energy to water in the direction of flow of water which is equal to 2000 J/m$^3$. P$_A$ and P$_B$ are pressure at section A and B respectively. Take $\rho$ (density of water) = 1000 kg/m$^3$. Mark the correct options.\par
\qfig{p2_phy_s1_q04_1}
\begin{qoptions}
\opt{A}{Speed of water at the outlet is 4 m/s.}
\opt{B}{Speed of water at the outlet is 1 m/s.}
\opt{C}{$P_A-P_B=1500\,\text{N/m}^2$}
\opt{D}{$P_A-P_B=5500\,\text{N/m}^2$}
\end{qoptions}
\qn{5} Charges are distributed in a spherical shell with an inner radius $R_1$ and an outer radius $R_2$. The charge density in the shell is $\rho=a+br$, where r is the distance from the center to the observation point. There is no charge distribution elsewhere in the space :-
\begin{qoptions}
\opt{A}{The electric field for $R_1<r<R_2$ is given by $E_2=\dfrac{1}{\varepsilon_0r^2}\left[\dfrac a3\left(r^3-R_1^3\right)+\dfrac b4\left(r^4-R_1^4\right)\right]$.}
\opt{B}{The electric potential for $R_1<r<R_2$ is zero.}
\opt{C}{The electric potential for $r<R_1$ is $V=\dfrac1{\varepsilon_0}\left[\dfrac a2\left(R_2^2-R_1^2\right)+\dfrac b3\left(R_2^3-R_1^3\right)\right]$.}
\opt{D}{The total charge in the spherical shell is $4\pi\left[\dfrac a3\left(R_2^3-R_1^3\right)+\dfrac b4\left(R_2^4-R_1^4\right)\right]$.}
\end{qoptions}
\qn{6} A circuit consists of two resistors $R_1$, $R_2$ and capacitor C as shown in figure. The capacitor consists of two circular plates with a radius b and separation d (b $\gg$ d). The key K is kept closed for a very long time. At t = 0, the key K is opened :- ($C=\dfrac{\epsilon_0\pi b^2}{d}$)\par
\qfig{p2_phy_s1_q06_1}
\begin{qoptions}
\opt{A}{The magnitude of displacement current just after the key is opened is $\dfrac{V_0}{R_1+R_2}$.}
\opt{B}{The magnitude of magnetic field for r $<$ b is $B=\dfrac{\mu_0r}{2\pi b^2}\cdot\dfrac{V_0}{R_1+R_2}e^{-\frac{1}{R_1C}t}$.}
\opt{C}{The magnitude of Poynting vector for r $<$ b is $S=\dfrac{R_1V_0^2r}{2\pi db^2(R_1+R_2)^2}e^{-\frac2{R_1C}t}$.}
\opt{D}{The magnitude of Poynting vector for r $<$ b is $S=\dfrac{R_2V_0^2r}{2\pi db^2(R_1+R_2)}e^{-\frac2{R_1C}t}$.}
\end{qoptions}
\sect{SECTION-I (ii)}
\paraq{7}{8}
A conductor of length L has the shape of a half cylinder of radius r. Thickness of the sheet used to make the conductor is t. The conductivity of the material varies with the angle$\theta$ as $\sigma=\sigma_0\cos\theta$. A battery of emf V and negligible internal resistance is connected across the two semi-circular ends.\par
\qfig{p2_phy_s2_para7_8_1}
\qn{7} Which of the following is true regarding the current density (j)?
\begin{qoptions}
\opt{A}{Current density is same at all point on the semi-circular cross-section.}
\opt{B}{Current density varies on the semi-circular cross-section, decreasing with $\theta$.}
\opt{C}{Current density varies on the semi-circular cross-section, increasing with $\theta$}
\opt{D}{Current density is same at every point in the conductor.}
\end{qoptions}
\qn{8} Current flowing through the battery is :-
\begin{qoptions}
\opt{A}{$\dfrac{\sigma_0Vrt}{2L}$}
\opt{B}{$\dfrac{\sigma_0Vrt}{L}$}
\opt{C}{$\dfrac{2\sigma_0VLt}{r}$}
\opt{D}{$\dfrac{2\sigma_0Vrt}{L}$}
\end{qoptions}
\paraq{9}{10}
In order to reduce the noise of the submarine and increase its forward speed, the propeller can be replaced by an electromagnetic propeller. There are two sets of left and right thrusters below the submarine. Each set is composed of 6 identical linear channel thrusters made of insulating materials. The principle diagram is as shown below. The linear channel is filled with seawater with a resistivity of $\rho$ = 0.2 ohm-m. In the space of a $\times$ b $\times$ c = 0.3 m $\times$ 0.4 m $\times$ 0.3 m in the channel, there is a uniform magnetic field generated by the superconducting coil, with a magnetic induction intensity B = 6.4 T and a direction perpendicular to the face of area a $\times$ c outward. On the magnetic field area, there are metal plates M and N with area (a $\times$ b = 0.3 m $\times$ 0.4 m). When they are connected to the dedicated DC power supply for the thruster, a wave from N to M is generated in the seawater between the two plates, and the current is 1000 A through each channel, assuming that this current only exists in the magnetic field area. Excluding the internal resistance of the power supply and the resistance of the wire, and taking the density of seawater $\rho_m$ = 1000 kg/m$^3$ :-\par
\qfig{p2_phy_s2_para9_10_1}
\qn{9} Mark the correct statement.\par
(1) Changing the magnitude of current in left propeller only can turn the submarine.\par
(2) Changing the direction of current in left propeller only (keeping magnitude constant) can turn the submarine.\par
(3) Changing the direction of current in both the propellers can reverse the submarine.\par
(4) Changing the magnitude of magnetic field in both the propellers can change the speed of submarine.
\begin{qoptions}
\opt{A}{1 only}
\opt{B}{1 and 2}
\opt{C}{2 and 3}
\opt{D}{1, 2, 3 and 4}
\end{qoptions}
\qn{10} When the submarine moves at a constant speed of 30 m/s, the power supplied by the DC source is :-
\begin{qoptions}
\opt{A}{$67\times10^5$ W}
\opt{B}{$5.5\times10^5$ W}
\opt{C}{$11.17\times10^5$ W}
\opt{D}{$56\times10^5$ W}
\end{qoptions}
\sect{SECTION-II (i)}
\paraq{1}{2}
An experimental setup of verification of photoelectric effect is shown in Figure. The voltage across the electrodes is measured with the help of an ideal voltmeter, and which can be varied by moving jockey `J' on the potentiometer wire. The battery used in potentiometer circuit is of 20 V and its internal resistance is 2 $\Omega$. The resistance of 100 cm long potentiometer wire is 8 $\Omega$.\par
\qfig{p2_phy_s3_para1_2_1}
The photocurrent is measured with the help of an ideal ammeter. Two plates of potassium oxide of area 50 cm$^2$ at separation 0.5 mm are used in the vacuum tube. Photocurrent in the circuit is very small, so we can treat the potentiometer circuit as an independent circuit. (Take : $\in_0$ = 8.85 $\times$ 10$^{-12}$ C$^2$/N-m$^2$)
\qn{1} The number of electrons that appear on the surface of the cathode plate are given as K $\times$ 10$^9$, when the jockey is connected at the end `P' of the potentiometer wire. Assume that no radiation is falling on the plates. Find the K.
\qn{2} When a light falls on the anode plate, the ammeter reading remains zero till jockey is moved from the end P to the middle point of the wire PQ. There after, the deflection is recorded in the ammeter. The maximum kinetic energy (in eV) of the emitted electron is :-
\paraq{3}{4}
As shown in the figure, the radii of the upper and lower plano-convex transparent cylinders are $R_1$ and $R_2$ respectively, their sections (planes) are parallel to their respective axes. The busbars passing through the tangent points are perpendicular to each other. Take the tangent point of the two cylinders as the origin of the rectangular coordinate system O-XYZ, the generatrix of the lower cylinder passing through the tangent point O is the X-axis, and the generatrix of the upper cylinder passing the tangent point O is the Y-axis. A beam of visible light with a wavelength in vacuum $\lambda$ is incident paraxially along the negative direction of the Z-axis, and the light reflected from the upper and lower cylinders will interfere near the origin and can be observed with the help of an optical reading microscope. The projection of the interference fringes on the upper cylinder onto the X-Y plane. $R_1$ and $R_2$ are much larger than the maximum gap between the two cylinders corresponding to the paraxial ray interference area. The refractive index of air is $n_0$ = 1.00.\par
\qfig{p2_phy_s3_para3_4_1}
\qn{3} The equation for 3$^{rd}$ bright fringe is $\dfrac{x^2}{R_1\lambda}+\dfrac{y^2}{R_2\lambda}=\dfrac\alpha2$. The value of $\alpha$ is :-
\qn{4} The equation for 4$^{th}$ dark fringe is $\dfrac{x^2}{R_1\lambda}+\dfrac{y^2}{R_2\lambda}=\dfrac\beta2$. The value of $\beta$ is :-
\paraq{5}{6}
Two identical and coaxial superconducting insulated circular coils 1 and 2 with radius R and self inductance L each are located far apart from each other. Each coil carries a current I in the same sense. Now both the coils are moved slowly so that both coils are in contact with each other.
\qn{5} The current through each coil in final configuration is $\alpha$I. The value of $\alpha$ is :-
\qn{6} The work done by external force in changing the configuration of the system is $\beta LI^2$. The value of $|\beta|$ is :-
\sect{SECTION-II (ii)}
\qn{7} A man of mass m starts moving w.r.t a platform of mass 2 m with a velocity u = 9/13 m/s towards right as shown in the figure. The platform is fitted with a concave mirror of focal length f. The velocity of image (in m/s) at the moment is :-\par
\qfig{p2_phy_s4_q07_1}
\qn{8} A person is observing a stone at the bottom of the pool. The water depth is h = 14 m. The angle between the line of sight and the normal line of the water surface is 60$^\circ$, the refractive index of water is 1.33. The depth of the image of the stone (in m) is :-
\qn{9} The half-life of the long lived isotope Radium-226 was originally measured by allowing it to charge an insulated metal plate and measuring the time taken to reach a known potential, repeatedly, over a long period of time. Such an apparatus is shown schematically in figure. The apparatus is supported on insulating rods and contained in a vacuum.\par
The Ra-226 is contained in the thin walled glass container in the form of 3.06 mg of newly prepared radium bromide (RaBr). As it beta decays, the metal plate that it sits on charges up and the flexible gold leaf rises until it reaches 220 V, when it touches the earthed electrode, it is discharged. Assume that all beta particles escape through the glass and away from the apparatus. Measurements on the apparatus show that the potential on the gold leaf is proportional  to the charge, with the constant of proportionality being 6.40 pC V$^{-1}$. The relative atomic masses of radium and bromine atoms are 226 and 80 respectively. The gold leaf rises and falls every 88 seconds. The half-life of Ra-226 is k $\times$ 10$^9$ sec. The value of k is: (Take : N$_A$ = 6 $\times$ 10$^{23}$ nuclei/mole and $\ell$n2 = 0.7)\par
\qfig{p2_phy_s4_q09_1}
\raggedcolumns
\partsec{2}{CHEMISTRY}
\sect{SECTION-I (i)}
\qn{1} Select the correct options
\begin{qoptions}
\opt{A}{Emulsion can be diluted with any amount of the dispersion medium}
\opt{B}{Emulsions can be broken into constituent liquid by heating, freezing, centrifuging etc.}
\opt{C}{Emulsifying agents for W/O emulsion are proteins, gums, natural and synthetic soaps etc}
\opt{D}{Butter and Cream are W/O type emulsions}
\end{qoptions}
\qn{2} Select the correct options
\begin{qoptions}
\opt{A}{Vapour pressure of a pure liquid is an intentive property and also a colligative property.}
\opt{B}{Solution showing very little positive deviation can form minimum boiling point azeotrope}
\opt{C}{Life time for zero order reaction is finite and zero order reaction are always complex in nature.}
\opt{D}{Electrolysis product of brine (aq. conc. NaCl) using inert electrodes are H$_2$ and Cl$_2$ gas at cathode and anode respectively.}
\end{qoptions}
\qn{3} Reaction scheme:\par
\qfig{p2_chem_s1_q03_1}
Among the following reactions, Identify the reaction(s) which gives Compound (Y) as a major product ?
\begin{qoptions}
\opt{A}{\qfigi{p2_chem_s1_q03_A}}
\opt{B}{\qfigi{p2_chem_s1_q03_B}}
\opt{C}{\qfigi{p2_chem_s1_q03_C}}
\opt{D}{\qfigi{p2_chem_s1_q03_D}}
\end{qoptions}
\qn{4} Some name of polymers and corresponding monomers are given below. Identify the correctly matched option(s)?
\begin{qoptions}
\opt{A}{Natural rubber: \qfigi{p2_chem_s1_q04_1}}
\opt{B}{Dacron: \qfigi{p2_chem_s1_q04_2}}
\opt{C}{Nylon-6: \qfigi{p2_chem_s1_q04_3}}
\opt{D}{Teflon: \qfigi{p2_chem_s1_q04_4}}
\end{qoptions}
\qn{5} How many among the following groups are of lewis acids
\begin{qoptions}
\opt{A}{BF$_3$, BCl$_3$, B$_2$H$_6$}
\opt{B}{H$^{+}$, NH$_4^{+}$, CH$_4$}
\opt{C}{BCl$_3$, AlCl$_3$, B$_2$H$_6$}
\opt{D}{CH$_3$CO$^{+}$, Fe$^{3+}$, Mg$^{2+}$}
\end{qoptions}
\qn{6} Correct statements about silicones is/are-
\begin{qoptions}
\opt{A}{These are organosilicon polymers}
\opt{B}{A famous example of silicones is zeolite}
\opt{C}{The chain length of the polymer can be controlled by adding (CH$_3$)$_3$ SiCl}
\opt{D}{Cu is used as catalyst in preparation of (CH$_3$)$_2$SiCl$_2$ from CH$_3$Cl + Si}
\end{qoptions}
\sect{SECTION-I (ii)}
\paraq{7}{8}
When $[\text{Co(NH}_3)_4(\text{H}_2\text{O})_2]^{3+}$ reacts with excess of NH$_3$, $[\text{Co(NH}_3)_5(\text{H}_2\text{O})]^{3+}$ is formed then HCl in added to form a complex (X) which contain no H$_2$O as a ligand \& shows purple colour
\qn{7} If one mole of (X) gives two moles of AgCl on treating with excess AgNO$_3$ solution then its formula and its magnetic nature is :-
\begin{qoptions}
\opt{A}{CoCl$_3$.6NH$_3$ , diamagnetic}
\opt{B}{CoCl$_3$.5NH$_3$ , diamagnetic}
\opt{C}{CoCl$_3$.4NH$_3$ , paramagnetic}
\opt{D}{CoCl$_3$.5NH$_3$ , paramagnetic}
\end{qoptions}
\qn{8} Select correct about this complex (X) :-
\begin{qoptions}
\opt{A}{It is an inner orbital octahedral complex.}
\opt{B}{Complex show geometrical isomerism}
\opt{C}{IUPAC name of complex is tetraaminedichloridocobalt(III) chloride}
\opt{D}{Van't Hoff factor of complex = 2}
\end{qoptions}
\paraq{9}{10}
Reaction scheme:\par
\qfig{p2_chem_s2_para9_10_1}
\qn{9} Identify Compound(S) and Reagent(R) in above reaction scheme?
\begin{qoptions}
\opt{A}{Compound(S) =\par\qfigi[0.75]{p2_chem_s2_q09_1}}
\opt{B}{Compound(S) = \qfigi{p2_chem_s2_q09_2}}
\opt{C}{Compound(S) = \qfigi{p2_chem_s2_q09_3}}
\opt{D}{Compound(S) = \qfigi{p2_chem_s2_q09_4}}
\end{qoptions}
\qn{10} Correct statement about compound (T) is?
\begin{qoptions}
\opt{A}{It is Chiral compound.}
\opt{B}{It has 3 Chiral center}
\opt{C}{It can be classified as Ketal}
\opt{D}{It can reduce fehling solution}
\end{qoptions}
\sect{SECTION-II (i)}
\paraq{1}{2}
An inorganic compound (A) reacts with dilute H$_2$SO$_4$ to give a colourless gas (B) which smells like rotten egg.\par
Black precipitate of (X) $\xleftarrow{\text{AgNO}_3}$ B $\xrightarrow{\text{Alk.Na}_2[\text{Fe(CN)}_5\text{NO}]}$ Violet coloured compound (Y)
\qn{1} Change in oxidation state of Fe, When Na$_2[\text{Fe(CN)}_5\text{NO}]$ reacts with B to give (Y).
\qn{2} Number of moles of acidified K$_2$Cr$_2$O$_7$ required to oxidise 340 gram gas B.\par
(Given atomic mass : H = 1, O = 16, C = 12, F = 19, N = 14, S = 32, P = 31, Cl = 35.5, Br = 80, I = 127, K = 39, Na = 23, Pb = 207, Al = 27, Si = 28, Mn = 55, Cr = 52)
\paraq{3}{4}
\qfig{p2_chem_s3_para3_4_1}
\qn{3} Total number of compounds from (P) To (U) which can give yellow ppt when react with NaOI?
\qn{4} Total number of Monochlorinated products obtained when compound(W) undergoes photochlorination?
\paraq{5}{6}
The unit cell is the smallest repeating unit in a crystal structure. The unit cell of gold crystal is found by X-ray diffraction to have the face-centred cubic unit structure. The side of the unit cell is found to be 0.400 nm. The density of Au is 2.0625 $\times$ 10$^4$ kg m$^{-3}$
\qn{5} If mass of one Au atom is X $\times$ 10$^{-25}$ kg. Then value of X will be
\qn{6} If the Avogadro constant is Y $\times$ 10$^{23}$, then value of Y is\par
(Given that the relative atomic mass of Au is 198.33)
\sect{SECTION-II (ii)}
\qn{7} $\Delta$G (in Joule) for given transformation\par
5 moles H$_2$O ($\ell$, 1bar, T K) $\rightarrow$ H$_2$O ($\ell$, 10 bar, T K)\par
(density of water = 1 gm/mL \& 1 lit - bar = 100 J)
\qn{8} EMF of the given galvanic cell at 25$^\circ$C\par
Pt$|$H$_2$(1bar)$|$RNH$_3^{+}$Cl$^{-}$ (1M) $||$ CH$_3$COOH (0.1 M)$|$H$_2$(1 bar)$|$Pt\par
is 0.12V, then pK$_b$ for RNH$_2$ (weak base) will be\par
Given that : K$_a$ (CH$_3$COOH) = 10$^{-5}$, and $\dfrac{2.303RT}{F}$ = 0.06.
\qfullbegin
\qn{9} For the given compound, the total number of optically active stereoisomers is?\par
\qfig{p2_chem_s4_q09_1}
\qfullsplit
\raggedcolumns
\qfullend
\partsec{3}{MATHEMATICS}
\sect{SECTION-I(i)}
\qn{1} Consider a $\Delta$ABC, whose vertices are A (3, $-$2), B (3, 1) \& C (1, 5). P is a point on altitude AD of $\Delta$ABC such that ordinate of P is 5. Now a circle is drawn taking P as centre, DE as one of the tangent, where E is the point on AB such that AB is perpendicular to CE. Identify the true statement.
\begin{qoptions}
\opt{A}{If tangents are drawn from point D \& E to given circle, then they intersect at a point on AC.}
\opt{B}{PDBE is a cyclic quadrilateral.}
\opt{C}{Centre of given circle is (17, 5)}
\opt{D}{$\angle$PBD = $\angle$DEP}
\end{qoptions}
\qn{2} Consider a circle $x^2+y^2=25$ with centre $P_1$. Pair to tangents are drawn from $P_2\left(-\dfrac{25}{3},0\right)$ to given circle which touch the circle at $T_1$ \& $T_2$. If 2 circles each of radius 3 are drawn with centre $C_1$ \& $C_2$ which touch the given circle externally at $T_1$ \& $T_2$ respectively, then
\begin{qoptions}
\opt{A}{Area of $\Delta P_1C_1C_2$ is $\dfrac{786}{25}$}
\opt{B}{Area of $\Delta P_2C_1C_2$ is $\dfrac{1969}{75}$}
\opt{C}{Area of $\Delta P_1C_1C_2$ is $\dfrac{768}{25}$}
\opt{D}{Area of $\Delta P_2C_1C_2$ is $\dfrac{1696}{75}$}
\end{qoptions}
\qn{3} Let the number of ways in which 3 positive integers a, b, c can be selected such that abc = 30030 is k, then
\begin{qoptions}
\opt{A}{If a, b, c be not necessarily distinct, then value of k is 726.}
\opt{B}{If a $<$ b $<$ c, then value of k is 121.}
\opt{C}{Value of k is 729 if a, b, c be not necessarily distinct.}
\opt{D}{If a $<$ b $<$ c, then value of k is 115.}
\end{qoptions}
\qn{4} Consider a parabola whose directrix passes through P(0, 9). If T(3, 2) is the foot of perpendicular drawn from focus S (2, $-$1) on a tangent of parabola, then $-$
\begin{qoptions}
\opt{A}{Length of latus rectum of parabola is $8\sqrt2$}
\opt{B}{Equation of tangent at vertex is x + y $-$ 5 = 0}
\opt{C}{Equation of axis of parabola is x $-$ y = 3}
\opt{D}{Directrix is at a distance $2\sqrt2$ from focus.}
\end{qoptions}
\qn{5} If $f(n,x)=\displaystyle\sum_{k=0}^{n}{}^nC_k\sin(kx)\cos((n-k)x)$, then
\begin{qoptions}
\opt{A}{$f\left(10,\dfrac\pi{20}\right)=512$}
\opt{B}{$f\left(5,\dfrac\pi{20}\right)$ satisfies the equation $x^2=256$}
\opt{C}{$f\left(11,\dfrac\pi{11}\right)=0$}
\opt{D}{$f\left(7,\dfrac\pi{14}\right)=256$}
\end{qoptions}
\qn{6} If $\left(\displaystyle\sum_{r=1}^{2016}2^{r-1}\tan2^{r-1}\theta\right)=\cot\dfrac\pi{2^\lambda}-2^\mu$ $\left(\text{where }\theta=\dfrac\pi{2^{2018}}\right)$, then identify correct option(s).
\begin{qoptions}
\opt{A}{$\lambda-\mu=2$}
\opt{B}{$\lambda+\mu=(\lambda-\mu)(\mu+1)$}
\opt{C}{$\lambda^\mu<\mu^\lambda$}
\opt{D}{$\lambda^\mu>\mu^\lambda$}
\end{qoptions}
\sect{SECTION-I(ii)}
\paraq{7}{8}
Consider the function $f(x)=\left(1+\dfrac1x\right)^x$, then
\qn{7} The function f (x)
\begin{qoptions}
\opt{A}{has a maxima but no minima}
\opt{B}{has a minima but no maxima}
\opt{C}{has exactly one maxima and one minima}
\opt{D}{has neither a maxima nor a minima}
\end{qoptions}
\qn{8} Range of the function f (x) is
\begin{qoptions}
\opt{A}{$(0,\infty)$}
\opt{B}{$(-\infty,e)$}
\opt{C}{$(1,\infty)$}
\opt{D}{$(1,e)\cup(e,\infty)$}
\end{qoptions}
\paraq{9}{10}
Suppose you do not know the function f (x), however some information about f (x) is listed below. Read the following carefully before attempting the questions.\par
(i) f (x) is continuous and defined for all real numbers\par
(ii) $f'(-5)=0$; $f'(2)$ is not defined and $f'(4)=0$\par
(iii) (-5, 12) is a point which lies on the graph of f(x)\par
(iv) $f''(2)$ is undefined, but $f''(x)$ is negative everywhere else.\par
(v) the signs of $f\,'$ are given below :\par
\qfig{p2_math_s2_para9_10_1}
\qn{9} On the possible graph of y= f(x) we have
\begin{qoptions}
\opt{A}{x = $-$5 is a point of relative minima}
\opt{B}{x = 2 is a point of relative maxima}
\opt{C}{x = 4 is a point of relative minima}
\opt{D}{graph of y = f (x) must have a geometrical sharp corner or cusp}
\end{qoptions}
\qn{10} From the possible graph of y = f(x), we can say that
\begin{qoptions}
\opt{A}{there is exactly one point of inflection on the curve}
\opt{B}{f(x) increases on $-5<x<2$ and x $>$4 and decreases on $-\infty<x<5$ and 2 $<$ x $<$ 4}
\opt{C}{the curve is always concave down}
\opt{D}{curve always concave up}
\end{qoptions}
\sect{SECTION-II (i)}
\paraq{1}{2}
Five balls are to be placed in three boxes. Boxes should hold all the five balls so that no box remains empty.
\qn{1} Find number of ways if balls and boxes are identical is
\qn{2} Find number of ways if balls as well as boxes are identical but boxes are kept in a row is
\paraq{3}{4}
Let S and T are two events defined on a sample space with probabilities P(S) = 0.5, P(T) = 0.69\par
$P\left(\dfrac ST\right)=0.5$
\qn{3} The value of twice of P (S and T)
\qn{4} The value of twice of P (S or T)
\paraq{5}{6}
Let \par\resizebox{\linewidth}{!}{$f(x)=\begin{cases}\displaystyle\lim_{n\to\infty}\left(\dfrac{ax}{n}\left(\sum_{K=1}^{a}\dfrac{\left[K^2-e^{-x}+K-1\right]}{K(K+1)}\right)\right)+\lambda, & x>0\\[2ex] b, & x=0\\[2ex] \displaystyle\lim_{n\to\infty}\sum_{K=1}^{a}\dfrac{\left\{K^2+K+e^{x}-1^2\right\}}{K(K+1)}, & x<0\end{cases}$}\par
(Where [$\cdot$] and \{$\cdot$\} denotes G.I.F and F.P. F respectively) is differentiable in R then
\qn{5} The value of a + b + $\lambda$ is equal to
\qn{6} The value of $f\,'(\log_e2)+f\,'(\log_e1/2)+f\,'(\log_e3/2^2)+\ f\,'(\log_e5/2^3)+\dots\infty$ is
\sect{SECTION-II (ii)}
\qn{7} Let A be a square matrix such that $A(\text{adj }A)=\begin{bmatrix}4&0&0\\0&4&0\\0&0&4\end{bmatrix}$ then $\dfrac{|\text{adj (adj A)}|}{|\text{adj A}|}$ is
\qn{8} If $z_1,z_2,z_3$ are three points lying on the circle $|z|=2$, then minimum value of $|z_1+z_2|^2+|z_2+z_3|^2+|z_3+z_1|^2$ is
\qn{9} If f(x) is a polynomial function satisfying.\par
$f(x)+f(y)+f(xy)-2=f(x)f(y)\forall x,y\in R$\par
and f(3) = 10, then $\dfrac{f(4)}{17}$ is equal to
\end{multicols}
\end{document}
